\section{Impacts environnementaux et sociétaux (DD\&RS)}

\subsection{Impacts liés à l’activité de l’entreprise}
L’activité RUHENNA est centrée sur l’artisanat et des prestations événementielles. Plusieurs éléments contribuent à une approche plus responsable :
\begin{itemize}
    \item \textbf{Produits naturels} : RUHENNA met en avant l’usage de produits naturels, avec notamment un approvisionnement auprès d’une entreprise française (\textbf{Aroma-Zone}) pour certains composants/matières premières.
    \item \textbf{Artisanat et petites séries} : la production en quantités maîtrisées peut limiter certains gaspillages (sur-stockage), même si elle nécessite une organisation rigoureuse.
\end{itemize}

\subsection{Impacts liés au numérique (plateforme web)}
La plateforme elle-même induit un impact numérique (hébergement, transfert de données, stockage). Les principaux postes sont :
\begin{itemize}
    \item \textbf{Images} : la galerie et la boutique impliquent beaucoup de médias ; c’est généralement la source principale de trafic et de poids des pages.
    \item \textbf{Hébergement} : déploiement sur Vercel et base de données PostgreSQL managée (Neon) ; l’impact dépend de l’usage réel (trafic, requêtes).
    \item \textbf{Transferts réseau} : consultations des pages, chargements d’images, interactions admin (API).
\end{itemize}

\subsection{Impacts liés à l’équipe / au contexte de travail}
Dans le contexte d’une petite structure, les impacts de l’équipe sont surtout indirects :
\begin{itemize}
    \item réduction de la charge administrative grâce à la centralisation (moins d’échanges redondants) ;
    \item meilleure planification des rendez-vous (optimisation des déplacements et de l’organisation).
\end{itemize}

\section{Actions de réduction des impacts}

\subsection{Actions mises en œuvre}
Plusieurs choix techniques ont été faits pour limiter l’impact, sans dégrader l’expérience utilisateur :
\begin{itemize}
    \item \textbf{Optimisation des images} : usage de \texttt{next/image} et stockage Cloudinary (CDN), pour servir des images adaptées aux tailles d’écran.
    \item \textbf{Architecture simple} : une application Next.js unique (UI + API) réduit la duplication d’infrastructures et la complexité de déploiement.
    \item \textbf{Modèle de données structuré} : PostgreSQL + Prisma limitent les incohérences et réduisent des traitements correctifs coûteux.
\end{itemize}

\noindent Le choix de \textbf{Cloudinary} s’explique par la place centrale des images (galerie, produits, blog) : plutôt que de gérer un stockage de fichiers “maison”, Cloudinary fournit un CDN et des mécanismes d’optimisation (formats plus légers, compression, tailles adaptées) qui réduisent le poids transféré vers les visiteurs \cite{cloudinaryDocs}. Couplé à \texttt{next/image}, cela aide à limiter les transferts inutiles tout en conservant une bonne qualité visuelle.

\subsection{Suppression des images côté Cloudinary}
Afin d’éviter de stocker inutilement des images non utilisées, la plateforme supprime aussi les médias côté Cloudinary lorsqu’un contenu est supprimé dans l’administration (galerie, bannière, produits, couverture d’article). Cette pratique limite le stockage “orphelin” et garde un système propre \cite{cloudinaryDocs}.